\documentclass[11pt,a4paper]{article}
\usepackage[margin=26mm]{geometry}
\usepackage{amsmath,amssymb,amsthm,hyperref}
\newtheorem{lemma}{Lemma}
\newtheorem{proposition}[lemma]{Proposition}
\newcommand{\Sym}{\operatorname{Sym}}
\newcommand{\Alt}{\operatorname{Alt}}
\title{Complete permutation searches and the six-point row}
\author{Human Compression project}\date{}
\begin{document}\maketitle
This note proves the complete search reduction for two experimental rows of
Tables~1 and~3 of Leemans's survey, and independently settles its six-point
symmetric row. The larger numerical conclusions require every listed branch
to finish and its complete output to be checked. A bounded or interrupted
computation gives no upper bound on the number of markings.

We use the previously proved enumeration theorem for actual permutation
markings. Its proof concerns a group $G$ normalized by a permutation group
$A$, rather than alternating groups alone. In the cases considered here,
$G=\Sym(n)$ or $\Alt(n)$ and $A=\Sym(n)$. For $n\ne6$ this action induces
the full automorphism group of $G$. Thus point conjugacy is exactly the
isomorphism relation on ordered markings in degrees fourteen and fifteen.
The exceptional six-point automorphism is handled explicitly below.

\section{The inherited recursion}
Let $I$ be the set of nonidentity involutions in $G$, sorted by the fixed
permutation ordering. A node carries a prefix
$t=(t_0,\ldots,t_{m-1})$, its actual generated subgroup $H$, its pointwise
stabilizer $S=C_A(t_0)\cap\cdots\cap C_A(t_{m-1})$, and the pool of
involutions commuting with $t_0,\ldots,t_{m-2}$. For each $S$-orbit of
candidates outside $H$, choose its least element. Append that element
exactly when the resulting tuple has the intersection property.
The new pool is obtained by intersecting the old pool with
$C_G(t_{m-1})$. Thus the newly appended mark need not commute with its
immediate predecessor, but every later mark must do so.

The intersection test is the acquired recursive criterion: the two end
omissions must themselves be string C-markings, and their actual
intersection must be the subgroup generated by the common middle marks.
The left omission is already certified at the preceding node; the test
recursively checks the right one and computes the actual intersection.
No quotient is used. A new generator lying in $H$ is forbidden by the
intersection property, so its rejection loses no valid extension.

For completeness, the induction on prefixes is short. Every extension
can first have its existing prefix conjugated to the representative
already retained. The remaining possible conjugacies fixing that prefix
are exactly $S$, so precisely one representative of each extension orbit
is visited. Conversely, two prefixes retained in different such orbits
cannot be conjugate with the preceding prefix fixed. Induction gives
both completeness and absence of repetitions.

The candidate set is invariant under $S$: every element of $S$ fixes
all selected marks, hence $H$, and preserves every commutation condition
imposed by the pool. Thus enumerating the orbits on this whole domain
is equivalent to closing the same set of seeds under the action.
The replacement of GAP's \texttt{Orbits} by \texttt{OrbitsDomain} changes
only that implementation choice. All six completed small branches have
identical full literal outputs under the two routines.

The next-pool evaluation has a further exact implementation simplification.
For a sorted pool $P\subseteq I$,
\[P\cap\{x\in I:[x,t]=1\}=\{x\in P:[x,t]=1\}.\]
Filtering $P$ directly by this predicate therefore gives exactly the same
ordered list as forming the complete ambient commuting set and intersecting
it with $P$. No representative, subgroup or descendant changes. All nine
completed comparison branches retain identical full literal output.

The transitivity rejection is also literal. Every further generator lies
in the current pool, which can only shrink. If the subgroup generated by
$H$ and that entire pool is intransitive, no descendant generates the
transitive target $G$. A node already generating $G$ has no independent
extension. Both rejections therefore preserve completeness. Transitive
proper subgroups are retained as well, but only records with actual group
order $|G|$ enter the present counts; as subgroups of $G$, those records
then generate $G$ itself.

\section{Connected prefixes for the fixed targets}
Suppose adjacent marks $t_j,t_{j+1}$ commute. The string relations then
make every mark in $\{t_0,\ldots,t_j\}$ commute with every mark in the
remaining interval. The two generated subgroups are nontrivial and have
trivial intersection by the intersection property. They therefore give
a nontrivial direct product decomposition of the whole group.

Neither target group admits such a decomposition for $n\ge5$.
For $A_n$ this follows from its standard simplicity. For $S_n$, any
nontrivial normal subgroup $N$ contains $A_n$: otherwise simplicity gives
$N\cap A_n=1$, so $[N,A_n]=1$. But $C_{S_n}(A_n)=1$, since a permutation
centralizing every three-cycle preserves every three-element support,
and hence fixes each point. Thus two nontrivial direct factors would
both contain $A_n$, contradicting their trivial intersection.
Consequently every adjacent pair in a target marking must fail to
commute. The search may reject a prefix at its first commuting adjacent
pair; later marks cannot restore a connected diagram across that cut.

This filter removes complete stabilizer orbits, because the stabilizer
fixes the preceding mark. All remaining minimum representatives are
unchanged. The new implementation enumerates the connected part of the
old transitive subgroup list; it is not a complete list of disconnected
proper subgroups. Complete small lists are compared after precisely this
filter, and all target counts must remain identical. The direct-factor
argument, rather than the small controls, proves the uniform target
completeness of the filter.

\section{Independent initial branches}
\begin{lemma}
For $G=\Sym(n)$, the nonidentity involution classes under $A=\Sym(n)$
are indexed by $1\leq e\leq\lfloor n/2\rfloor$, where $e$ is the number
of transpositions. For $G=\Alt(n)$, retain exactly the even values of $e$.
Let $u_e$ be the first member of $I$ having $e$ transpositions. The entire
branch beginning with $u_e$ is obtained by starting the inherited
recursion at
\[
       ((u_e),\langle u_e\rangle,C_A(u_e),I).
\]
\end{lemma}
\begin{proof}
Two involutions are conjugate in the symmetric group precisely when their
cycle types agree, and the cycle type is $2^e1^{n-2e}$. Such an involution
is even precisely when $e$ is even. The sorted first member is the least
element of this orbit, exactly the representative chosen at the empty
prefix. The inherited recursion imposes no commutation restriction on
the second mark, so the child pool is still all of $I$. Its subgroup and
stabilizer are exactly the displayed groups. Every subsequent operation
is unchanged. Different values of $e$ cannot produce conjugate ordered
markings because their first entries have different cycle types.
\end{proof}

The direct entry removes only the computation of all initial orbits.
It does not change candidate selection or the tests below that entry.
The original whole search, its version selecting one initial orbit after
forming them all, and the direct entry have been compared on $S_5,S_7$
and $A_5$. Their complete literal lists agree after taking the appropriate
union. The known counts provide additional controls. These checks support
the implementation; the lemma supplies the all-degree justification.

Each initial branch has a separate immutable driver and a bounded
execution receipt. Completed branches may be reused without rerunning
them. A cap records incompleteness. Numerical admission requires one
complete receipt for every admissible value of $e$, with the same group,
rank range and engine edition. A separately documented equivalent entry
edition may be combined only through the lemma and exact provenance;
an unfinished predecessor is never counted as a second branch.

\section{Rank ranges and reversal}
For $S_{14}$, the only unacquired ranks are $3\leq r\leq8$.
The accepted stability theorem supplies ranks $9$ through $13$, since
with $k=14-r$ its condition $14\geq2k+3$ holds. The accepted independent
generating-length bound excludes ranks at least $14$.
For $A_{15}$, the accepted alternating maximal-rank theorem gives
$r\leq7$. Hence the fixed searches require seven initial branches for
$S_{14}$ through rank eight and three initial branches for $A_{15}$
through rank seven.

Let $c_r$ be the resulting number of ordered classes and $s_r$ the number
fixed by reversal. For each literal representative the latter condition
is tested by asking whether a permutation in $\Sym(n)$ conjugates its
entire ordered tuple to its reverse. The number of classes with dual
markings identified is
\[
                         (c_r+s_r)/2.
\]
This follows from the orbits of the reversal involution on the complete
set of ordered classes. Reversal can exchange distinct initial branches;
therefore this quotient is formed only after all branches are combined.
The two large target groups have no exceptional outer automorphism affecting
this calculation.


\section{The complete six-point calculation}
\begin{proposition}
For $S_6$, the numbers of string C-markings of ranks $3,4,5$, up to group
isomorphism and reversal, are respectively $2,4,1$.
\end{proposition}
\begin{proof}
All three initial involution classes are exhausted through rank five.
The recursion gives $6,14,2$ ordered classes under point conjugacy.
The literal representatives are included with every subgroup and output
receipt. In addition to the recursive enumeration test, a separate verifier
checks every pair of subsets of each representative's index set, requiring
that the intersection of the two actual generated subgroups equals the
subgroup on the common indices. It also checks generation of $S_6$ and
all involution and string relations. There are $6016$ such subset pairs.

Point conjugacy is not the full equivalence relation in degree six.
We use the standard fact that $\operatorname{Aut}(S_6)$ has order $1440$,
with its inner subgroup of order $720$. The certificate constructs the full
automorphism group and an actual bijection $\phi$ taking a transposition to
three disjoint transpositions. It verifies that the inner automorphisms and
$\phi$ generate the full group. For every retained tuple, conjugacy under
$S_6$ locates the unique retained class of its $\phi$-image and of its reverse.
These give two literal permutations on each complete class list. Their
orbits, not a point-conjugacy formula borrowed from another degree, give the
required numbers:
\[
\begin{array}{c|rrr}
\text{rank}&3&4&5\\\hline
\text{point-conjugacy classes}&6&14&2\\
\text{full automorphism classes}&3&7&1\\
\text{full automorphisms and reversal}&2&4&1
\end{array}
\]
The exact two permutations are retained in the certificate. Inner
conjugacies act trivially on the class list, so the group generated by these
two permutations is precisely the required equivalence relation.
The accepted independent generating-length bound gives rank at most five
for $S_6$, which completes this row.
\end{proof}

The number four in the middle column is the previously unacquired rank-four
cell of the survey. The other two values are fresh checks of acquired cells.
The input does not use these expected counts: complete lists are generated
first, and source comparison follows. The six-point controls also test the
connected-prefix change on proper transitive subgroups: exactly seventeen
disconnected markings are removed across the three initial classes, while
all twenty-two target markings remain unchanged.

The search and its proof are reused from the acquired finite-census
method. The direct initial entry, invariant-domain evaluation, direct-factor
pruning, the pool evaluation and exact survey-cell interfaces are the additional work
described here. The source's numerical answers do
not direct the search and are compared only with completed results.
This component supplies no claim about an open conjecture, no new
measurement, and no whole-source admission.
\end{document}
